\documentclass[10pt,a4paper]{amsart}
\usepackage[margin=19mm]{geometry}
\usepackage{amsmath,amssymb,mathtools}
\usepackage[scr=rsfs,cal=euler]{mathalfa}
\usepackage{xfrac}
\usepackage[unicode,hidelinks,bookmarks=false]{hyperref}
\newcommand{\ie}{i.e.\ }
\newcommand{\cf}{cf.\ }
\newcommand{\resp}{respectively\ }
\newcommand{\Subsection}{\subsection}
\providecommand{\nobreakdash}{\nobreak\hbox{-}}
\setlength{\emergencystretch}{2em}
\multlinegap0pt
\usepackage{tikz}
\usetikzlibrary{cd}
\usepackage{mathtools}
\usepackage{amssymb}
\usepackage{float}
\usepackage{xcolor}
\usepackage{bm}
\usepackage{tikz}
\usepackage[shortlabels]{enumitem}
\newtheorem{lemma}{Lemma}
\newtheorem{proposition}{Proposition}
\usepackage{cleveref}
\crefname{lemma}{Lemma}{Lemmas}
\crefname{proposition}{Proposition}{Propositions}
\def\C{\mathbb{C}}
\newcommand{\M}{\mathrm{M}}
\def\P{\mathbb{P}}
\def\Q{\mathbb{Q}}
\def\Z{\mathbb{Z}}
\newcommand{\abs}[1]{\left\lvert#1\right\rvert}
\newcommand{\cL}{\mathcal{L}}
\newcommand{\cO}{\mathcal{O}}
\def\calA{\mathcal{A}}
\def\calL{\mathcal{L}}
\newcommand{\cl}{\textrm{cl}}
\renewcommand{\emptyset}{\varnothing}
\title{Magic positivity of Snapper polynomials \\ for matroids}
\author{Shiyue Li}
\date{Source: arXiv:2609.24917v1 (CC BY 4.0)}
\begin{document}
\maketitle

\noindent\emph{Source and licence.} Magic positivity of Snapper polynomials \\ for matroids. Version arXiv:2609.24917v1 (CC BY 4.0), distributed by arXiv under the Creative Commons Attribution 4.0 International licence, http://creativecommons.org/licenses/by/4.0/, as recorded in the arXiv licence metadata for this exact version and shown on the abstract page https://arxiv.org/abs/2609.24917v1. Full licence notice: \texttt{licenses/arxiv-2609.24917-CC-BY-4.0.txt}.\par\smallskip

\noindent\textit{Editorial scope.} Counted unit: the article's proposition prop:pullback (source0.tex lines 1641--1651). The article's complete original proof of it is delivered with it (source0.tex lines 1653--1696, one \texttt{proof} environment of 44 lines). No mathematics is written, repaired or re-derived: the statement and the proof are the article's own lines, in the article's own notation, and no retained line is substituted or re-flowed.  Retained uncounted as the source-local context the proof needs: lines 927--931; lines 951--957; lines 1627--1630.  Nothing was omitted from any retained span, so the package carries no omission record.  No retained line contains a citation, so the wrapper carries no bibliography and no bibliography entry.  The retained lines contain the article's own diagram code, which the wrapper typesets with the standard tikz packages; no image file is delivered and the package declares no asset dependency.  Editorial formatting only, in addition: an AMS \texttt{amsart} preamble that restates the article's own declarations and macros used by the retained lines, namely the environments \texttt{lemma}, \texttt{proposition} on separate counters, together with the standard packages those lines need.  The retained environments print in this document as Lemma 1, Proposition 1 in order of appearance, whereas the article prints the same items with the numbering of its own full document.  The complete native source of the article as distributed in the arXiv e-print for this exact version is bundled unchanged as \texttt{sources/211239/source0.tex}.
\begin{equation}\label{eq:LP-flats}
  \cL_P \;\coloneqq\; \bigotimes_{F} \cL_F^{\otimes \omega_{\M, P}(F)} \;\in\; K(\M),
  \qquad
  \chi_{\M,P}(q) \;\coloneqq\; \chi\big(\M, \cL_P^{\otimes q}\big) \in \Q[q].
\end{equation}

\begin{lemma}
  \label{lem:support-general}
  Let $\M$ be a loopless matroid, and $P = \sum_{S} y_S\,\Delta_S$ with $y_S \in \Z$. For every flat $F$ of rank at least $2$,
  \[
    \omega_{\M, P}(F) \;=\; \sum_{\cl_{\M}(S) = F} y_S .
  \]
\end{lemma}

\begin{equation}\label{eq:OP}
  \cO_{X_{n-1}}(P) \;=\; \bigotimes_{\emptyset \ne S \subseteq E}
  \cL_{\Delta_S}^{\otimes y_S} .
\end{equation}

\begin{proposition}\label{prop:pullback}
  Let $\M$ be loopless and realized by an arrangement $\calA \subseteq L$.  For
  every nonempty $S \subseteq E$,
  \[
    \iota^{*} \cL_{\Delta_S} \;\cong\; \cL_{\cl_{\M}(S)} .
  \]
  As a consequence, for every lattice generalized permutohedron $P$,
  \[
    \iota^{*}\, \cO_{X_{n-1}}(P) \;\cong\; \cL_P .
  \]
\end{proposition}

\begin{proof}
  For the first statement, the polytope $\Delta_S$ is the polytope of the toric variety $\P^{|S|-1}$,
  and, since $\Sigma_{n-1}$ refines its normal fan, it defines a toric morphism
  \[
    p_{S} \colon X_{n-1} \longrightarrow \P^{|S|-1}
  \]
  such that $\cL_{\Delta_S} \cong p_S^{\ast}\cO(1)$.  Let $T^{n-1}$ denote the dense torus
  of $X_{n-1}$, which can be identified with the dense torus of $\P^{n-1}$.  On the dense torus $T^{n-1}$, the map $p_S$ is the coordinate projection
  \[
    [x_e]_{e \in E} \longmapsto [x_e]_{e \in S} .
  \]

  By the definition of the closure of $S$ in the matroid $\M$,
  $\bigcap_{e \in S} H_e = \bigcap_{e \in \cl_{\M}(S)} H_e$, and
  $L^{S} = L^{\cl_{\M}(S)}$.  The coordinates $x_e$ with $e \in S$ embed $L^{S}$
  into $\C^{S}$, realizing $\P(L^{S})$ as a linear subspace of $\P^{|S|-1}$.
  Therefore, the aforementioned maps form a commutative square
  \[
    \begin{tikzcd}[row sep=2.2em, column sep=3em]
      W_{\calA} \arrow[r, hook, "\iota"] \arrow[d, "\pi_{\cl_{\M}(S)}"'] & X_{n-1} \arrow[d, "p_S"] \\
      \P(L^{\cl_{\M}(S)}) \arrow[r, hook] & \P^{|S|-1}.
    \end{tikzcd}
  \]
  Pulling back $\cO(1)$ on $\P^{\abs{S}-1}$ along the two paths
  now gives the first statement, and since $\cO(1)$ on $\P^{\abs{S}-1}$ restricts to $\cO(1)$ on $\P(L^{\cl_{\M}(S)})$, we have
  \[
    \iota^{*}\cL_{\Delta_S} \;\cong\; (p_S \circ \iota)^{*}\,\cO(1)
    \;=\; \pi_{\cl_{\M}(S)}^{*}\,\cO(1) \;=\; \cL_{\cl_{\M}(S)} .
  \]

  For the second statement, expand $\cO_{X_{n-1}}(P)$ by \eqref{eq:OP}, apply the
  previous statement to each factor, and group the subsets by their closure:
  \[
    \iota^{*}\, \cO_{X_{n-1}}(P)
    \;\cong\;
    \bigotimes_{\emptyset \ne S \subseteq E} \cL_{\cl_{\M}(S)}^{\otimes y_S}
    \;=\;
    \bigotimes_{\emptyset \ne F \in \calL(\M)}
      \cL_{F}^{\otimes\, \sum_{\cl_{\M}(S) = F}\, y_S} .
  \]
  For flats of rank at least two, \Cref{lem:support-general} identifies the
  exponent with $\omega_{\M,P}(F)$; the remaining factors are trivial. Thus
  the product is $\cL_P$ by \eqref{eq:LP-flats}.
\end{proof}

\end{document}
